\documentclass[10pt,a4paper]{amsart}
\usepackage[margin=19mm]{geometry}
\usepackage{amsmath,amssymb,mathtools}
\usepackage[scr=rsfs,cal=euler]{mathalfa}
\usepackage{xfrac}
\usepackage[unicode,hidelinks,bookmarks=false]{hyperref}
\newcommand{\ie}{i.e.\ }
\newcommand{\cf}{cf.\ }
\newcommand{\resp}{respectively\ }
\newcommand{\Subsection}{\subsection}
\providecommand{\nobreakdash}{\nobreak\hbox{-}}
\setlength{\emergencystretch}{2em}
\multlinegap0pt
\usepackage{amssymb}
\usepackage{mathtools}
\newtheorem{proposition}{Proposition}
\title{Normal Ordering and Stirling-Type Combinatorics for Double Ore Extensions of Type $(14641)$}
\author{Andr\'es Rubiano}
\date{Source: arXiv:2606.08396v1 (CC BY 4.0)}
\begin{document}
\maketitle

\noindent\emph{Source and licence.} Normal Ordering and Stirling-Type Combinatorics for Double Ore Extensions of Type $(14641)$. Version arXiv:2606.08396v1 (CC BY 4.0), distributed by arXiv under the Creative Commons Attribution 4.0 International licence, http://creativecommons.org/licenses/by/4.0/, as recorded in the arXiv licence metadata for this exact version and shown on the abstract page https://arxiv.org/abs/2606.08396v1. Full licence notice: \texttt{licenses/arxiv-2606.08396-CC-BY-4.0.txt}.\par\smallskip

\noindent\textit{Editorial scope.} Counted unit: the article's proposition PropositionInternalNormalOrdering (source0.tex lines 615--635). The article's complete original proof of it is delivered with it (source0.tex lines 637--686, one \texttt{proof} environment of 50 lines). No mathematics is written, repaired or re-derived: the statement and the proof are the article's own lines, in the article's own notation, and no retained line is substituted or re-flowed.  No further source-local context is needed: every cross-reference of every retained line resolves inside the two delivered spans.  Nothing was omitted from any retained span, so the package carries no omission record.  No retained line contains a citation, so the wrapper carries no bibliography and no bibliography entry.  No retained line contains a figure environment, an image-inclusion command or drawing code, so no image is delivered and the package declares no asset dependency.  Editorial formatting only, in addition: an AMS \texttt{amsart} preamble that restates the article's own declarations and macros used by the retained lines, namely the environments \texttt{proposition} on separate counters, together with the standard packages those lines need.  The retained environments print in this document as Proposition 1 in order of appearance, whereas the article prints the same items with the numbering of its own full document.  The complete native source of the article as distributed in the arXiv e-print for this exact version is bundled unchanged as \texttt{sources/202328/source0.tex}.
\begin{proposition}\label{PropositionInternalNormalOrdering}
Let $q,r,p,s\in\Bbbk$ with $q,p\in\Bbbk^\times$. Assume that
\[
x_2x_1=qx_1x_2+rx_1^2,
\qquad
y_2y_1=py_1y_2+sy_1^2.
\]
Then, for all $m,n\in\bar{\mathbb{N}}$, the following identities hold:
\begin{align}
x_2^nx_1^m
&=
\sum_{k=0}^{n}
\Lambda_{q,r}(n,m;k)\,
x_1^{m+k}x_2^{n-k}, \label{EquationXInternalNormalOrdering}\\
y_2^ny_1^m
&=
\sum_{k=0}^{n}
\Lambda_{p,s}(n,m;k)\,
y_1^{m+k}y_2^{n-k}. \label{EquationYInternalNormalOrdering}
\end{align}
\end{proposition}

\begin{proof}
We prove \eqref{EquationXInternalNormalOrdering}; the proof of \eqref{EquationYInternalNormalOrdering} is identical after replacing $(x_1,x_2,q,r)$ by $(y_1,y_2,p,s)$.

First observe that
\[
x_2x_1^m
=
q^mx_1^mx_2+r[m]_qx_1^{m+1}.
\]
This follows by induction on $m$. The case $m=1$ is the defining relation. If the identity holds for $m$, then
\begin{align*}
x_2x_1^{m+1}
&=
\left(q^mx_1^mx_2+r[m]_qx_1^{m+1}\right)x_1\\
&=
q^mx_1^m(qx_1x_2+rx_1^2)+r[m]_qx_1^{m+2}\\
&=
q^{m+1}x_1^{m+1}x_2+r(q^m+[m]_q)x_1^{m+2}\\
&=
q^{m+1}x_1^{m+1}x_2+r[m+1]_qx_1^{m+2}.
\end{align*}

Now write
\[
x_2^nx_1^m
=
\sum_{k=0}^{n}
C_{n,m}(k)x_1^{m+k}x_2^{n-k}.
\]
For $n=0$ one has $C_{0,m}(0)=1$. Multiplying the expansion for $n$ on the left by $x_2$ and using the previous identity gives
\[
C_{n+1,m}(k)
=
q^{m+k}C_{n,m}(k)
+
r[m+k-1]_qC_{n,m}(k-1),
\]
where $C_{n,m}(-1)=0$. A direct verification shows that
\[
C_{n,m}(k)
=
r^kq^{m(n-k)}
\begin{bmatrix}
n\\
k
\end{bmatrix}_q
[m]_{q}^{\overline{k}}
\]
satisfies this recurrence together with the initial condition.
\end{proof}

\end{document}
